\scanpage{265}\mspage{181}
% archive sheet number: 264
est toujours ``meilleure'' que celle-ci (si celle-ci
\deleted{\unsure{lui}} n'est \deleted{\unsure{pas déjà}} déjà canonique, i.e. si les
deux ne sont égales), à cause des deux
faits suivants : a) \inserted{\unsure{surtout|d'abord}} parce qu'elle est \uline{canonique},
i.e. ne dépend d'aucun autre choix que
celui du seul cardinal $\pi_0$, et b) parce
qu'il y a a priori plus de cardinaux
utiles et de cardinaux réguliers pour
la fcg canonique. \inserted{Donc au total, il me semble que les filtrations cardinales \unsure{g.} les seules utiles sont celles qui sont \deleted{\unsure{access}} canoniques.}
Notons qu'une filtration cardinale (au sens
de SGA 4 \unsure{I \S|I}, que je conserve ici) \deleted{qui} est de plus
\uline{accessible} \inserted{(i.e. $\mathcal F^\pi\unsure{\subset|=}M_\pi$ pour tt $\pi\geq\pi_0$)}, \unsure{n'est|en est} autre qu'une \deleted{\unsure{filtr}}
fcga canonique, telle que \uline{tous} les
cardinaux sont réguliers (a fortiori utiles)
pour \unsure{ladite|la dite} filtration. Je présume \deleted{que}
(sans avoir cependant construit d'exemple)
que les filtrations cardinales canoniques
des catégories de la forme $M=\operatorname{Ind}_{\pi_0}(C)$,
$C$ une \inserted{petite} catégorie \uline{karoubienne} (pour
que $\pi_0$ soit un cardinal utile) \deleted{\unsure{don}}
ne sont pas forcément des fc (non généralisées, i.e. au sens SGA), i.e. je présume
qu'il arrive qu'il y ait des cardinaux $\pi\geq\pi_0$
non réguliers, et $=$ superflus.
